\section{Diagnóstico do modelo selecionado}

O modelo \sarimafinal foi reestimado na amostra completa após a seleção com base no teste. O diagnóstico residual avaliou centralização dos resíduos, autocorrelação residual, distribuição aproximada e comportamento temporal dos erros.

O modelo assume que, após uma diferenciação regular e uma diferenciação sazonal de período 12, a série se torna aproximadamente estacionária. Na escala diferenciada, a dependência temporal é representada pelos componentes autorregressivos e de médias móveis especificados, com erros de média zero, variância aproximadamente constante e ausência de autocorrelação serial. A especificação é univariada, não inclui variáveis exógenas nem tendência determinística explícita. Essas suposições são avaliadas por meio do diagnóstico dos resíduos.

\begin{figure}[H]
\centering
\includegraphics[width=0.92\textwidth]{diagnostico_residuos_modelo_final.png}
\caption{Diagnóstico dos resíduos do modelo selecionado}
\label{fig:diagnostico}
\end{figure}

O ponto crítico do ajuste está no teste de Ljung--Box. Os p-valores foram baixos nas defasagens 6, 12, 18 e 24, indicando rejeição da hipótese de ausência de autocorrelação residual. A Tabela \ref{tab:ljung} resume os resultados.

\begin{table}[H]
\centering
\caption{Teste de Ljung--Box para os resíduos}
\label{tab:ljung}
\begin{tabular}{rrr}
\toprule
Defasagem & Estatística & p-valor \\
\midrule
6 & 23,36 & 0,0007 \\
12 & 77,82 & $1{,}07 \times 10^{-11}$ \\
18 & 79,97 & $8{,}68 \times 10^{-10}$ \\
24 & 80,92 & $4{,}34 \times 10^{-8}$ \\
\bottomrule
\end{tabular}
\end{table}

Portanto, o modelo selecionado apresenta bom desempenho preditivo no teste, mas não elimina completamente a dependência serial dos resíduos. Essa limitação sugere que a série pode conter choques, quebras estruturais ou componentes não capturados por um SARIMA simples. Como o enunciado solicita um relatório curto e focado em ARIMA/SARIMA, essa evidência foi registrada como limitação, em vez de expandir o trabalho para modelos mais complexos.

